% Options for packages loaded elsewhere
\PassOptionsToPackage{unicode}{hyperref}
\PassOptionsToPackage{hyphens}{url}
\PassOptionsToPackage{dvipsnames,svgnames,x11names}{xcolor}
\documentclass[
  10pt,
  fontset=fandol]{ctexart}
\usepackage{xcolor}
\usepackage[margin=16mm]{geometry}
\usepackage{amsmath,amssymb}
\setcounter{secnumdepth}{-\maxdimen} % remove section numbering
\usepackage{iftex}
\ifPDFTeX
  \usepackage[T1]{fontenc}
  \usepackage[utf8]{inputenc}
  \usepackage{textcomp} % provide euro and other symbols
\else % if luatex or xetex
  \usepackage{unicode-math} % this also loads fontspec
  \defaultfontfeatures{Scale=MatchLowercase}
  \defaultfontfeatures[\rmfamily]{Ligatures=TeX,Scale=1}
\fi
\usepackage{lmodern}
\ifPDFTeX\else
  % xetex/luatex font selection
\fi
% Use upquote if available, for straight quotes in verbatim environments
\IfFileExists{upquote.sty}{\usepackage{upquote}}{}
\IfFileExists{microtype.sty}{% use microtype if available
  \usepackage[]{microtype}
  \UseMicrotypeSet[protrusion]{basicmath} % disable protrusion for tt fonts
}{}
\makeatletter
\@ifundefined{KOMAClassName}{% if non-KOMA class
  \IfFileExists{parskip.sty}{%
    \usepackage{parskip}
  }{% else
    \setlength{\parindent}{0pt}
    \setlength{\parskip}{6pt plus 2pt minus 1pt}}
}{% if KOMA class
  \KOMAoptions{parskip=half}}
\makeatother
\setlength{\emergencystretch}{3em} % prevent overfull lines
\providecommand{\tightlist}{%
  \setlength{\itemsep}{0pt}\setlength{\parskip}{0pt}}
\usepackage{amsmath,amssymb,mathtools,needspace}
\xeCJKDeclareCharClass{CJK}{"2032,"0400->"04FF,"2160->"216B,"2460->"2473,"25A0,"25C7,"25CB,"25CF}
\IfFontExistsTF{Arial Unicode MS}{\xeCJKsetup{AutoFallBack=true}\setCJKfallbackfamilyfont{\CJKrmdefault}{Arial Unicode MS}}{}
\setcounter{secnumdepth}{0}
\setlength{\emergencystretch}{2em}
\allowdisplaybreaks
\usepackage{bookmark}
\IfFileExists{xurl.sty}{\usepackage{xurl}}{} % add URL line breaks if available
\urlstyle{same}
\hypersetup{
  pdftitle={用正交函数作最小二乘拟合},
  colorlinks=true,
  linkcolor={Maroon},
  filecolor={Maroon},
  citecolor={Blue},
  urlcolor={Blue},
  pdfcreator={LaTeX via pandoc}}

\title{用正交函数作最小二乘拟合}
\author{}
\date{}

\begin{document}
\maketitle

\subsubsection{3.6.2 用正交函数作最小二乘拟合}\label{ux7528ux6b63ux4ea4ux51fdux6570ux4f5cux6700ux5c0fux4e8cux4e58ux62dfux5408}

用最小二乘法得到的法方程组 (3.6.6)，其系数矩阵 \(G\) 是病态的，但如果 \(\varphi_0(x),\varphi_1(x),\cdots,\varphi_n(x)\) 是关于点集 \(\{x_i\}\)（\(i=0,1,\cdots,m\)）带权 \(\omega(x_i)\)（\(i=0,1,\cdots,m\)）正交的函数族，即

\[
(\varphi_j,\varphi_k)=\sum_{i=0}^m\omega(x_i)\varphi_j(x_i)\varphi_k(x_i)=\begin{cases}0,&j\ne k,\\ A_k>0,&j=k,\end{cases}\tag{3.6.8}
\]

则方程组 (3.6.6) 的解为

\[
a_k^*=\frac{(f,\varphi_k)}{(\varphi_k,\varphi_k)}=\frac{\displaystyle\sum_{i=0}^m\omega(x_i)f(x_i)\varphi_k(x_i)}{\displaystyle\sum_{i=0}^m\omega(x_i)\varphi_k^2(x_i)}\quad(k=0,1,\cdots,n),\tag{3.6.9}
\]

且平方误差为

\[
\|\delta\|_2^2=\|f\|_2^2-\sum_{k=0}^n A_k(a_k^*)^2.
\]

现在根据给定节点 \(x_0,x_1,\cdots,x_m\) 及权函数 \(\omega(x)>0\)，造出带权 \(\omega(x)\) 正交的多项式 \(\{P_n(x)\}\)。注意 \(n\leqslant m\)，用递推公式表示 \(P_k(x)\)，即

\[
\begin{cases}
P_0(x)=1,\\
P_1(x)=(x-\alpha_1)P_0(x),\\
P_{k+1}(x)=(x-\alpha_{k+1})P_k(x)-\beta_kP_{k-1}(x)\quad(k=1,2,\cdots,n-1).
\end{cases}\tag{3.6.10}
\]

这里 \(P_k(x)\) 是首项系数为 \(1\) 的 \(k\) 次多项式。根据 \(P_k(x)\) 的正交性，得

\[
\left\{\begin{aligned}
a_{k+1}&=\frac{\displaystyle\sum_{i=0}^m\omega(x_i)x_iP_k^2(x_i)}{\displaystyle\sum_{i=0}^m\omega(x_i)P_k^2(x_i)}=\frac{(xP_k(x),P_k(x))}{(P_k(x),P_k(x))}=\frac{(xP_k,P_k)}{(P_k,P_k)},\\
\beta_k&=\frac{\displaystyle\sum_{i=0}^m\omega(x_i)P_k^2(x_i)}{\displaystyle\sum_{i=0}^m\omega(x_i)P_{k-1}^2(x_i)}=\frac{(P_k,P_k)}{(P_{k-1},P_{k-1})}
\end{aligned}\right.\quad(k=0,1,\cdots,n-1).\tag{3.6.11}
\]

\begin{quote}
校注（agent 补充）：式 (3.6.11) 第一行原书印为拉丁字母 \(a_{k+1}\)，式 (3.6.10) 则为希腊字母 \(\alpha_{k+1}\)，已分别保留。该式末的共同范围原书为 \(k=0,1,\cdots,n-1\)；若直接用于第二行 \(\beta_k\)，\(k=0\) 会涉及未定义的 \(P_{-1}\)。依式 (3.6.10)，实际递推使用的是 \(\beta_1,\ldots,\beta_{n-1}\)。见 \href{../assets/ch03b/p082-recursion-coefficients.png}{原式裁切}。
\end{quote}

\href{../../source-images/pdf-083.jpeg}{查看原页}

下面用归纳法证明这样给出的 \(\{P_k(x)\}\) 是正交的。由式 (3.6.10) 第二式及式 (3.6.11) 中 \(a_1\) 的表达式，有

\[
\begin{aligned}
(P_l,P_1)&=(P_0,xP_0)-\alpha_1(P_0,P_0)\\
&=(P_0,xP_0)-\frac{(xP_0,P_0)}{(P_0,P_0)}(P_0,P_0)=0.
\end{aligned}
\]

现假定 \((P_1,P_s)=0\)（\(l\ne s\)）对于 \(s=0,1,\cdots,l-1\) 及 \(l=0,1,\cdots,k\)（\(k<n\)）均成立，要证 \((P_{k+1},P_s)=0\) 对于 \(s=0,1,\cdots,k\) 均成立。由式 (3.6.10)，有

\[
\begin{aligned}
(P_{k+1},P_s)&=((x-\alpha_{k+1})P_k,P_s)-\beta_k(P_{k-1},P_s)\\
&=(xP_k,P_s)-\alpha_{k+1}(P_k,P_s)-\beta_k(P_{k-1},P_s),
\end{aligned}\tag{3.6.12}
\]

由归纳法假定，当 \(0\leqslant s\leqslant k-2\) 时

\[
(P_k,P_s)=0,\quad(P_{k-1},P_s)=0.
\]

另外，\(xP_s(x)\) 是首项系数为 \(1\) 的 \(s+1\) 次多项式，它可用 \(P_0,P_1,\cdots,P_{s+1}\) 的线性组合表示，而 \(s+1\leqslant k-1\)，故由归纳法假定又有

\[
(xP_k,P_s)\equiv(P_k,xP_s)=0,
\]

于是由式 (3.6.12)，当 \(s\leqslant k-2\) 时，\((P_{k+1},P_s)=0\)。

再看

\[
(P_{k+1},P_{k-1})=(xP_k,P_{k-1})-\alpha_{k+1}(P_k,P_{k-1})-\beta_k(P_{k-1},P_{k-1}),\tag{3.6.13}
\]

由假定有

\[
(P_k,P_{k-1})=0,
\]

\[
(xP_k,P_{k-1})=(P_k,xP_{k-1})=\left(P_k,P_k+\sum_{j=0}^{k-1}C_jP_j\right)=(P_k,P_k).
\]

利用式 (3.6.11) 的 \(\beta_k\) 表达式及以上结果，得

\[
(P_{k+1},P_{k-1})=(xP_k,P_{k-1})-\beta_k(P_{k-1},P_{k-1})=(P_k,P_k)-(P_k,P_k)=0.
\]

最后，由式 (3.6.11) 有

\[
\begin{aligned}
(P_{k+1},P_k)&=(xP_k,P_k)-\alpha_{k+1}(P_k,P_k)-\beta_k(P_k,P_{k-1})\\
&=(xP_k,P_k)-\frac{(xP_k,P_k)}{(P_k,P_k)}(P_k,P_k)=0.
\end{aligned}
\]

至此已证明了由式 (3.6.10) 及式 (3.6.11) 确定的多项式，\(\{P_k(x)\}\)（\(k=0,1,\cdots,n;n\leqslant m\)）组成一个关于点集 \(\{x_i\}\) 的正交系。

用正交多项式 \(\{P_k(x)\}\) 的线性组合作最小二乘曲线拟合，只要在根据式 (3.6.10) 及式 (3.6.11) 逐步求 \(P_k(x)\) 的同时，相应计算出系数

\[
a_k=\frac{(f,P_k)}{(P_k,P_k)}=\frac{\displaystyle\sum_{i=0}^m\omega(x_i)f(x_i)P_k(x_i)}{\displaystyle\sum_{i=0}^m\omega(x_i)P_k^2(x_i)}\quad(k=0,1,\cdots,n),
\]

并逐步把 \(\alpha_k^*P_k(x)\) 累加到 \(F(x)\) 中去，最后就可得到所求的拟合曲线

\[
y=F(x)=\alpha_0^*P_0(x)+\alpha_1^*P_1(x)+\cdots+\alpha_n^*P_n(x),
\]

\begin{quote}
校注（agent 补充）：本页第一式左端原扫描为 \((P_l,P_1)\)（第一下标为字母 \(l\)），依右端计算应对应 \((P_0,P_1)\)。下一段归纳假设原印为 \((P_1,P_s)=0\)，括号内条件却为 \(l\ne s\)，疑应为 \((P_l,P_s)=0\)。均忠实保留并见 \href{../assets/ch03b/p083-induction-start.png}{原文裁切}。页尾系数公式用 \(a_k\)，其后累加与拟合曲线改用 \(\alpha_k^*\)，也按原印保留；此处应是同一组拟合系数的记号不一致，不应与递推系数 \(\alpha_k\) 混用。见 \href{../assets/ch03b/p083-fit-coefficients.png}{原文裁切}。
\end{quote}

\href{../../source-images/pdf-084.jpeg}{查看原页}

这里 \(n\) 可事先给定或在计算过程中根据误差确定。

用这种方法编程序不用解方程组，只用递推公式，并且当逼近次数增加一次时，只要把程序中循环数加 \(1\)，其余不用改变。这是目前用多项式作曲线拟合最好的计算方法，有通用的语言程序供用户使用。

\begin{center}\rule{0.5\linewidth}{0.5pt}\end{center}

\subsection{相关原页校注（agent 补充）}\label{ux76f8ux5173ux539fux9875ux6821ux6ce8agent-ux8865ux5145}

以下校注来自本节覆盖的原页，可能也涉及同页相邻小节；原文照录与校正说明分开保留。

\subsubsection{原书 PDF 页 84 的校注}\label{ux539fux4e66-pdf-ux9875-84-ux7684ux6821ux6ce8}

\href{../pages/pdf-084.md}{完整原页转录} · \href{../../source-images/pdf-084.jpeg}{原页图像}

校注记录：ch03b-E012, ch03b-E013。

\begin{quote}
校注（agent 补充）：3.6.3 原书的 \(S_n\) 求和从 \(k=1\) 开始，但后续使用 \(F(a_0,a_1,\ldots,a_n)\) 和 \(a_k\;(k=0,\ldots,n)\)；段中系数列表印为 \(a_0,a_2,\ldots,a_n\)；内积公式右端第二因子的原印下标为 \(i\)，而左端为 \(j\)，相容写法应使用 \(\varphi_j\)。这些原书记号不一致均已保留，见 \href{../assets/ch03b/p084-multivariate-formulas.png}{原文裁切}。
\end{quote}

\subsection{本节覆盖原页的校注记录}\label{ux672cux8282ux8986ux76d6ux539fux9875ux7684ux6821ux6ce8ux8bb0ux5f55}

下列记录按原页关联，含同页相邻小节的校注；计算时同时读取上文校注和完整原页。

\begin{itemize}
\tightlist
\item
  \texttt{ch03b-E008}：原书 PDF 页 82
\item
  \texttt{ch03b-E009}：原书 PDF 页 82
\item
  \texttt{ch03b-E010}：原书 PDF 页 83
\item
  \texttt{ch03b-E011}：原书 PDF 页 83
\item
  \texttt{ch03b-E012}：原书 PDF 页 84
\item
  \texttt{ch03b-E013}：原书 PDF 页 84
\end{itemize}

\end{document}
